\section{From the functional inequalities to mixing}\label{sec:dynamics}

We now turn the two functional endpoints into worst-start mixing
bounds. Both deductions use the unscaled heat-bath form and keep the
factor $n$ between the two clocks explicit. The support-profile proof
is the finite-state truncation mechanism of
\cite[Theorem~1.1 and Lemma~2.1]{GMT2006}. The classical logarithmic
Sobolev proof uses the hypercontractive smoothing of
\cite[Theorems~3.5 and~3.7]{DSC1996}, followed by spectral-gap decay.

\subsection{The support-profile route}
\label{prof:mixing-transfer}

\begin{lemma}\label{lem:profile-to-mixing}
Let $\mu$ be a full-support law on the cube, with heat-bath form
$\D$, generator $\mathcal L$, and one-attempt kernel $P$. Suppose
that for some $\alpha>0$ and every nonnegative $f$ with support mass
$0<p=\mu(f>0)\le1/2$,
\begin{equation}
 \D(f)\ge \alpha\,\mu(f^2)\log(1/p).
 \label{eq:common-support-profile}
\end{equation}
Write $t_{\mathrm{mix}}^{\mathrm{ct}}(\varepsilon)$ and
$t_{\mathrm{mix}}^{\mathrm{dt}}(\varepsilon)$ for the worst-start
total-variation mixing times of $e^{t\mathcal L}$ and $P^k$, respectively.
Put
\[
 \mu_{\min}=\min_x\mu(x),\qquad
 M=\max\{8,\mu_{\min}^{-1}\},\qquad
 k_0=\left\lceil\frac{2n}\alpha
            \log\frac{\log(M/4)}{\log2}\right\rceil.
\]
Then, for every $0<\varepsilon<1/2$,
\begin{align}
 t_{\mathrm{mix}}^{\mathrm{ct}}(\varepsilon)
 &\le\frac1\alpha\log\frac{\log(M/4)}{\log2}
      +\frac1{2\alpha\log2}\log\frac7{4\varepsilon^2},
 \label{prof:continuous-mixing}\\
 t_{\mathrm{mix}}^{\mathrm{dt}}(\varepsilon)
 &\le k_0+
    \left\lceil\frac n{\alpha\log2}
                       \log\frac7{4\varepsilon^2}\right\rceil.
 \label{prof:discrete-mixing}
\end{align}
In particular, if $\log(1/\mu_{\min})\le Cn$, then for every fixed
$\varepsilon>0$ the worst-start mixing time is
$O_{C,\varepsilon}(\alpha^{-1}\log(n+1))$ in per-site time and
$O_{C,\varepsilon}(n\alpha^{-1}\log(n+1))$ attempts.
\end{lemma}

\begin{proof}
First obtain a Poincare inequality, which will handle the final part
of mixing. Subtract a median from a real function and take the positive
and negative parts $f_+,f_-$. Both supports have mass at most $1/2$.
For each edge of the cube, splitting into these two parts can only
decrease the sum of squared increments. The two-point representation
of each conditional variance therefore gives
$\D(f_+)+\D(f_-)\le\D(f)$. Applying
\eqref{eq:common-support-profile} to the two parts yields
\[
 \Var_\mu(f)\le\mu[(f-\operatorname{med}_\mu f)^2]
 \le\frac{\D(f)}{\alpha\log2}.
\]
Thus the gap is at least $\alpha\log2$.

If $g\ge0$ is a probability density relative to $\mu$ and $V=\mu(g^2)$,
let $u=(g-V/4)_+$. Markov's inequality and the elementary inequality
$(x-a)_+^2\ge x^2-2ax$ for $x,a\ge0$ give
\[
 \mu(u\ne0)\le4/V,\qquad \mu(u^2)\ge V/2.
\]
The truncation is a contraction, so $\D(g)\ge\D(u)$.
For $V\ge8$, \eqref{eq:common-support-profile} consequently implies
\begin{equation}
 \D(g)\ge\frac\alpha2 V\log(V/4)
       \ge\frac{\alpha V}{6}\log V.
 \label{prof:density-energy}
\end{equation}

For the continuous-time semigroup with generator
$\mathcal L=\sum_{i=1}^n(Q_i-I)$, where $Q_i$ is conditional expectation given
$X_{-i}$, the density from a point evolves by $g_t=e^{t\mathcal L}g_0$.
Reversibility gives $V'_t=-2\D(g_t)$. While $V_t\ge8$, setting
$z_t=\log(V_t/4)$ in \eqref{prof:density-energy} yields
$z'_t\le-\alpha z_t$. Put
\[
 t_0=\alpha^{-1}\log\frac{\log(M/4)}{\log2}.
\]
Since $V_0=1/\mu(x)\le M$, by time $t_0$ we have $V_t\le8$ for every
starting point.

Afterwards the gap bound gives
$V_{t_0+t}-1\le7e^{-2\alpha(\log2)t}$.
The inequality $\norm{g\mu-\mu}_{\TV}\le\tfrac12\sqrt{V-1}$
therefore proves \eqref{prof:continuous-mixing}.

For attempted updates we obtain the squared-norm decrease directly
from positivity of the average of conditional-expectation projections.
The transition operator is $P=n^{-1}\sum_iQ_i=I+\mathcal L/n$.
Each $Q_i$ is an orthogonal projection in $L^2(\mu)$, hence
$0\preceq P\preceq I$ and $P^2\preceq P$.
If $g_{k+1}=Pg_k$, then
\[
 V_k-V_{k+1}
 =\langle g_k,(I-P^2)g_k\rangle_\mu
 \ge\langle g_k,(I-P)g_k\rangle_\mu
 =\frac1n\D(g_k).
\]
While $V_k\ge8$, this and \eqref{prof:density-energy} give
$V_{k+1}\le V_k[1-\alpha\log(V_k/4)/(2n)]$.

If the next value is still at least $8$, taking logarithms gives
\[
 z_{k+1}\le(1-\alpha/(2n))z_k
          \le e^{-\alpha/(2n)}z_k,
 \qquad z_k=\log(V_k/4).
\]
The factor is positive: the gap bound and $-\mathcal L\preceq nI$ imply
$\alpha\log2\le n$. In addition, whenever this recurrence is used its
bracket must be positive since it bounds the positive value $V_{k+1}$.
Thus after at most $k_0$ steps the squared density norm is at most $8$.

On centered functions the same positive-operator inequality and the gap
imply
$V_{k+1}-1\le(1-\alpha\log2/n)(V_k-1)$.
The total-variation inequality above now gives
\eqref{prof:discrete-mixing}. This proves the transfer for deterministic
numbers of attempts.

If $\log(1/\mu_{\min})\le Cn$, then
$\log(\log(M/4)/\log2)=O_C(\log(n+1))$.
The two explicit bounds give the stated big-$O$ consequences for fixed
$0<\varepsilon<1/2$. For $\varepsilon\ge1/2$, the bound at accuracy $1/4$
suffices by monotonicity of mixing time in the accuracy threshold.
\end{proof}

The support bound in \Cref{prof:thm-profile} holds simultaneously
for all tests on its good-disorder event, so it applies throughout the
density evolution in the lemma. Fixed constants and logarithmic losses
are absorbed by choosing a slightly smaller exponent tolerance first.
For the big-$O$ consequence used here, the remaining hypothesis is a lower
bound on point masses. On
$\norm J\le3$, the Hamiltonian oscillation is at most $3n$, and
\begin{equation}\label{eq:minimum-mass}
 \mu_{J,\min}\ge\exp[-(3+\log2)n].
\end{equation}
This proves the asserted worst-start upper exponents through the
support-profile route.

\subsection{The logarithmic Sobolev route}

The classical inequality gives a second, independent smoothing
argument. With our convention
$\Ent_\mu(f^2)\le C_{\mathrm{LS}}\D(f)$,
finite-chain hypercontractivity takes the form
\begin{equation}\label{eq:hypercontractivity}
 \norm{e^{t\mathcal L}f}_{q(t)}\le\norm f_p,
 \qquad q(t)-1=(p-1)e^{4t/C_{\mathrm{LS}}},\quad p>1.
\end{equation}
Here the factor four is the reversible-chain normalization in
\cite[Theorem~3.5(ii)]{DSC1996}: its energy-to-entropy constant is
$1/C_{\mathrm{LS}}$ for our unscaled form.
For the density of a point mass, choose
$p=1+1/\max(1,\log(1/\mu_{\min}))$. Its $L^p$ norm is at most $e$.
The time at which $q(t)=2$ is
$O(C_{\mathrm{LS}}\log(2+\log(1/\mu_{\min})))$.
Expansion of the logarithmic Sobolev inequality at $f=1+u g$, with
$\mu g=0$, gives
\begin{equation}\label{eq:lsi-poincare}
 2\Var_\mu(g)\le C_{\mathrm{LS}}\D(g).
\end{equation}
The remaining $L^2$ decay to any fixed TV accuracy therefore costs
$O_\varepsilon(C_{\mathrm{LS}})$.

To transfer this bound to attempted updates, compare the two
semigroups in $L^2$. If $\theta\in[0,1]$ is an eigenvalue of $P$, then
$\theta^{2k}\le e^{-2k(1-\theta)}$. Expanding a centered initial
density in an orthonormal eigenbasis gives
\begin{equation}\label{eq:discrete-l2-comparison}
 \norm{P^k(g-1)}_2
 \le\norm{e^{(k/n)\mathcal L}(g-1)}_2.
\end{equation}
Thus the same argument, with a factor $n$ in time, proves the discrete
upper bound. Using \Cref{cap:thm-lsi} and \eqref{eq:minimum-mass}
establishes both upper exponents by the second route as well.

\subsection{Sharp exponents and preservation of the lower quantifiers}

The linear tests already give lower bounds for the full functional
constants: by definition $R_J\ge R_{\mathrm{lin},J}$.
Applying \Cref{thm:accepted-critical} and \eqref{eq:lsi-poincare} gives
\[
 C_{\mathrm{LS},J}\ge2R_J\ge2R_{\mathrm{lin},J}
                  =n^{2/3+o(1)}
\]
as a lower-exponent statement in disorder probability. The median
consequence of \Cref{prof:support-profile} gives the upper bound for
$R_J$, while \Cref{cap:thm-lsi} gives the upper bound for
$C_{\mathrm{LS},J}$ (and, by \eqref{eq:lsi-poincare}, another upper bound
for $R_J$). Intersecting these finitely many fixed-tolerance disorder
events gives the functional assertions of \Cref{thm:main}.

For a lower mixing bound at every fixed threshold below $1/2$,
choose a nonconstant eigenfunction $g$ of $-\mathcal L_J$ with
eigenvalue $1/R_J$. Center it and normalize $\norm g_\infty=1$.
At a point where $|g|=1$, its expectation decays exactly as
$e^{-t/R_J}$, whereas its stationary expectation is zero. Testing TV
against this bounded function gives
\[
 \max_x\norm{e^{t\mathcal L_J}(x,\cdot)-\mu_J}_{\TV}
 \ge\tfrac12 e^{-t/R_J}.
\]
The corresponding eigenvalue of $P_J$ is $1-1/(nR_J)$, so the discrete
lower bound is $\tfrac12(1-1/(nR_J))^k$.
For every fixed $\varepsilon<1/2$, these inequalities give mixing
lower bounds of order $R_J$ and $nR_J$, respectively. Together with
the upper bounds, they prove \Cref{thm:main}.

The realized-start assertions in \Cref{thm:accepted-critical} apply to
every deterministic little-oh time sequence. Its covariance and
linear-Rayleigh assertions apply to every diverging deterministic
sequence $M_n$. These quantifiers are retained from the companion;
no stronger upper quantifier or cutoff statement is inferred from
the new exponent estimates.

\subsection{Uniform relative-entropy decay}

The spectral-gap and classical logarithmic Sobolev bounds together
determine the critical time scale for uniform relative-entropy decay.
We state the normalization
explicitly, since constants called modified logarithmic Sobolev constants
use several conventions. For the per-site generator, define
\[
 \mathcal E_J(u,v):=-\mu_J[u\mathcal L_Jv]
 =\sum_i\mu_J\!\left[\Cov_{\mu_J}(u,v\mid X_{-i})\right],
 \qquad \mathcal E_J(f,f)=\D_J(f),
\]
and let
\[
 C_{\mathrm{EP},J}:=
 \sup_{\substack{g>0,\ \mu_Jg=1\\ \Ent_{\mu_J}(g)>0}}
 \frac{\Ent_{\mu_J}(g)}{\mathcal E_J(g,\log g)}.
\]
This is the inverse constant for the logarithmic entropy-production
inequality. The standard comparison and semigroup interpretation appear
in \cite[Propositions~3.5--3.6 and Theorem~2.4]{BobkovTetali2006Modified};
we give the short derivation in the present normalization.

\begin{corollary}[Entropy-production exponent]\label{cor:entropy-production}
For every realization of the disorder,
\[
 \frac{R_J}{2}\le C_{\mathrm{EP},J}\le\frac{C_{\mathrm{LS},J}}4.
\]
Consequently, for the critical zero-field Gaussian model,
$C_{\mathrm{EP},J}=n^{2/3+o(1)}$ in disorder probability, with the fixed
exponent-slack meaning of \Cref{thm:main}. For every initial probability
law $\nu$ on the cube and every $t\ge0$,
\[
 \operatorname{KL}(\nu e^{t\mathcal L_J}\Vert\mu_J)
 \le e^{-t/C_{\mathrm{EP},J}}\operatorname{KL}(\nu\Vert\mu_J)
 \le e^{-4t/C_{\mathrm{LS},J}}\operatorname{KL}(\nu\Vert\mu_J).
\]
\end{corollary}

\begin{proof}
For positive $a,b$,
$(a-b)(\log a-\log b)\ge4(\sqrt a-\sqrt b)^2$.
Apply this inequality to the two-point conditional representation of
each term of $\mathcal E_J$. The classical inequality then gives
\[
 \mathcal E_J(g,\log g)\ge4\D_J(\sqrt g)
 \ge\frac4{C_{\mathrm{LS},J}}\Ent_{\mu_J}(g),
\]
which proves the upper comparison. For a centered real function $f$,
$g=1+uf$ is a positive density for sufficiently small $|u|$. On the
finite cube,
\[
 \Ent_{\mu_J}(1+uf)=\tfrac12u^2\Var_{\mu_J}(f)+O(u^3),\qquad
 \mathcal E_J(1+uf,\log(1+uf))=u^2\D_J(f)+O(u^3).
\]
Letting $u\to0$ and then taking the supremum over $f$ proves
$C_{\mathrm{EP},J}\ge R_J/2$. For a requested exponent slack $\delta$,
apply \Cref{thm:main} with slack $\delta/2$ and then absorb the fixed
factors for sufficiently large $n$.

If $g_t$ is the density of $\nu e^{t\mathcal L_J}$ relative to $\mu_J$,
reversibility gives $g_t=e^{t\mathcal L_J}g_0$ and
\[
 \frac{d}{dt}\Ent_{\mu_J}(g_t)
 =-\mathcal E_J(g_t,\log g_t).
\]
The chain is irreducible, so $g_t>0$ for $t>0$ even when $g_0$ vanishes
at some states. The defining inequality and Gronwall give the displayed
decay; continuity includes $t=0$. Conversely, differentiating any
uniform exponential entropy bound at zero for positive initial
densities recovers its entropy-production inequality. Thus
$1/C_{\mathrm{EP},J}$ is the optimal rate $c$ in bounds of the displayed
form $\operatorname{KL}(\nu e^{t\mathcal L_J}\Vert\mu_J)
\le e^{-ct}\operatorname{KL}(\nu\Vert\mu_J)$ that hold for every initial
law and every time.
\end{proof}

If $\operatorname{KL}(\nu\Vert\mu_J)\le H$ for a fixed $H$, Pinsker's
inequality therefore gives total-variation accuracy $\varepsilon>0$ in
$O_{H,\varepsilon}(C_{\mathrm{LS},J})$ per-site time. A point mass has initial
relative entropy $\log(1/\mu_J(x))=O(n)$ under
\eqref{eq:minimum-mass}, recovering the logarithmic factor in the
worst-start upper bound. The exponent for $C_{\mathrm{EP},J}$ concerns
this uniform rate; it does not assert that every initial law relaxes at
the slowest rate.
